\subsection*{Definiciones operativas de medición}
Los umbrales anteriores conservan las Bases. Las definiciones siguientes instrumentan la medición propuesta: cada contacto, ticket y transacción conserva su identificador y las marcas de tiempo de origen. Los relojes contractuales y las reglas que requieran acuerdo del CLIENTE se distinguen de la instrumentación del registro (Distribuidora Puelche S.A., 2026a, arts. 78--79, p. 40).

La espera de un contacto es la diferencia entre recepción y primera atención. El indicador de respuesta de mesa divide contactos atendidos antes de veinte segundos por contactos recibidos y multiplica por cien; el abandono divide abandonos por contactos recibidos. La resolución al primer contacto divide primeros contactos resueltos sin escalamiento por primeros contactos atendidos. Una nueva llamada sobre un caso conserva el ticket y constituye otro contacto cuando efectivamente vuelve a ingresar a la cola. El registro permite auditar esa distinción y la calibración con abandono de T-15.

La respuesta al incidente mide desde recepción hasta primera actuación de atención registrada; el acuse automático del ticket se conserva como evento distinto. La restauración se registra cuando vuelve la función y se verifica la transacción. El cierre conforme registra la confirmación de la persona usuaria. Escalamiento, relevo, espera de tercero y espera fuera de cobertura conservan sus marcas y la hora de recepción original. Se informa el tiempo corrido de respuesta y restauración y se identifica la franja de atención. Para media se etiqueta la cobertura hábil extendida de la mesa conforme al Caso; para baja se registra el calendario hábil acordado por sitio y canal. Mientras el detalle no esté formalizado, el plazo se reporta en tiempo corrido con la franja correspondiente y sin descontar espera de terceros, relevo o tiempo fuera de cobertura (B-06).

El MTTR crítico es la suma de tiempos de restauración de incidentes críticos del mes dividida por su cantidad. No se sustituye por el tiempo hasta el cierre administrativo.

La tasa de cambios fallidos divide despliegues fallidos por despliegues de Producción del mes, con el fallo enlazado al cambio. La satisfacción es el promedio de respuestas válidas en la escala de uno a cinco y muestra cuántas respuestas la sustentan. Si un denominador es cero se informa «sin base de cálculo» con el conteo, sin asignar un 100\,\% ficticio. Los eventos con igual causa raíz se vinculan al mismo problema y conservan su fecha para verificar el máximo trimestral.

El RTO se mide desde la interrupción hasta la recuperación funcional comprobada. Para pérdida de región o de un sitio con extracción activa, el RPO de los servicios críticos compara el instante de pérdida con el último dato consistente recuperable fuera del dominio afectado. La pérdida de enlaces sin pérdida física comprueba autonomía y posterior reconciliación. Si al aislamiento total le sigue pérdida física, se mide desde esta última contra el último dato consistente externo y se conserva también el inicio del aislamiento. Los ensayos registran escenario, dominio, canal y evidencia de operaciones; B-11 identifica el control aún por acreditar (BA, art. 78.3; BTT RT-07.04; SD4, 4.3.2.4 y AL-DR-01).

\subsection*{Recorridos y ventana de observación}
Cada recorrido conserva el identificador de la transacción, canal, sitio, instante, resultado y correlación con el incidente. El siguiente catálogo define la comprobación inicial de los recorridos críticos y los p95 del caso. Su calendario por servicio y sitio se registra con el CLIENTE, incluidos los días hábiles, feriados, turnos y franjas excepcionales. Cada ficha conserva sitio, canal, horario local con zona America/Santiago, vigencia, responsables y referencia de aprobación. La ventana de 05:30--07:00 de lunes a sábado se controla de manera independiente.

\Needspace{0.48\textheight}
\begin{tablalafrox}{Comprobación de recorridos de negocio}{tab:sd10-recorridos}{F{0.28} F{0.30} Y}{\cab{Recorrido} & \cab{Ventana de referencia} & \cab{Resultado verificable}}
Preparación e identidad en bodega & Preparación 22:00--06:00 & Sesión autorizada y línea confirmada con persistencia local y marca de sincronización \\
Despacho y bloqueo térmico & Despacho 05:30--07:00, lunes a sábado & Validación térmica y documental de la carga. Si cambia, guía nueva del ERP antes de liberarla \\
Preventa, stock y crédito & Jornada del proceso registrada por canal & Línea aceptada y consulta de stock/crédito respondida; estado local y de integración identificables \\
Entrega & Jornada de ruta registrada por canal & Entrega registrada con su evidencia y estado de sincronización \\
\end{tablalafrox}
\fuentesdiez{BA, arts. 78--79, p. 40. Caso, RT-09.01 y RT-10.05, pp. 26--27. SD4, apartados 4.2.4.3 y 4.3.2.5.}

Para cada recorrido, sitio y mes se registra $T_0$, el tiempo de su calendario aplicable, y $A$, la duración conjunta de las ventanas programadas y aprobadas que intersectan ese calendario. La metodología propuesta usa $T=T_0-A$. El numerador $I$ suma una vez los intervalos sin resultado válido, dejando fuera sólo su intersección con dichas ventanas. Los excesos computan. El reporte conserva $I$, $T_0$, $A$ y $T$ para reproducir ambos resultados mientras se formaliza el denominador. Si éste es cero se informa sin base de cálculo. Los p95 se calculan por operación y se reportan aparte fallos y exclusiones.

Cada intervalo conserva el inicio del impacto, la primera recuperación funcional comprobada, la evidencia del canal y la referencia al ticket. Si falta observación, se informa su alcance y se contrasta con las trazas locales y la medición del CLIENTE. Se conserva el impacto parcial por recorrido, sede o canal. El despacho verifica además su umbral de indisponibilidad cero (BA, arts. 78--79).

\subsection*{Frescura analítica, avisos y telemetría}
El responsable de Operación utiliza los objetivos de la Tabla~\ref{tab:sd10-frescura}. NOC registra la primera desviación con su dominio, fuente y hora, y escala a Datos o a Calidad según el recorrido. Las metas de aviso remoto se informan por camino y tramo. Se identifican por separado los eventos sin red móvil ni cielo y la duración de su atraso, con protección local activa.

\Needspace{0.48\textheight}
\begin{tablalafrox}{Objetivos analíticos y térmicos del SD4}{tab:sd10-frescura}{F{0.23} F{0.25} Y}{\cab{Señal} & \cab{Objetivo o disparador} & \cab{Medición y evidencia}}
Analítica & Día $\leq$ 5 min; cierre $\leq$ 2 h; gestión $\leq$ 4 h & Desde transacción hasta disponibilidad analítica según vista. Datos verifica evento, carga y tablero. \\
Aviso móvil a Calidad & $\leq$ 2 min en al menos 99\,\% & Detección hasta aviso recibido, en eventos con camino móvil. M9 conserva ruta, ensayo, entrega y acuse. \\
Aviso satelital a Calidad & $\leq$ 5 min en al menos 95\,\% & Misma medición con cielo y vía satelital. M9 y Calidad contrastan tramo y condición de ruta en V2/V18. \\
Latido local del kit & Envío cada 10 s; vigencia menor de 30 s & TC58e comprueba firma, edad y consistencia antes de cada nueva bajada. Conserva alarma y decisión local. \\
Kit sin reportar & 6 min de silencio de ambos caminos & M9 registra la última recepción, abre alerta y añade llamada a Calidad. Se verifica acuse y actuación por separado. \\
\end{tablalafrox}
\fuentesdiez{BTT RT-05.29. SD4, capa analítica del apartado 4.1.3, PDF pp. 41 y 49; módulo M9, pp. 53--54; anexos 4-A, evento Kit Sin Reportar, y 4-W, V2, V6, V7 y V18.}

La falta de latido o su contradicción impiden habilitar una nueva bajada, conforme a GU-06. La pérdida durante una entrega ya habilitada se registra y exige latido válido antes de la siguiente. Calidad conserva la autoridad sobre la disposición sanitaria.

\subsection*{Conservación de evidencia de observabilidad}
La Tabla~\ref{tab:sd10-retencion} reúne los plazos declarados por SD4. Son clases distintas: una traza técnica no sustituye un registro de auditoría ni un documento sanitario o tributario. SD5 mantiene la retención por dominio y los controles de eliminación.

\Needspace{0.55\textheight}
\begin{tablalafrox}{Conservación declarada de observabilidad}{tab:sd10-retencion}{F{0.26} F{0.29} Y}{\cab{Clase} & \cab{Conservación} & \cab{Archivo y control}}
Registros técnicos & 12 meses en línea & 24 meses adicionales en archivo; 36 meses en total \\
Métricas & 13 meses en CloudWatch Metrics & Sin compromiso adicional de archivo identificado en la fuente \\
Trazas & 30 días en línea en \texttt{aws/spans} & Sin archivo adicional de trazas; política de Logs declarada en SD4 \\
Seguridad y auditoría & 7 años & Protección inmutable mediante Object Lock; retención de su propia clase \\
Telemetría local pendiente & Almacenamiento temporal de 24 h en ADOT & Exportar al restituir el enlace; comprobar capacidad con carga de diseño \\
\end{tablalafrox}
\fuentesdiez{SD4, apartados 4.1.3.7--4.1.3.8 y 4.2.4.5; SD5, apartado 5.2.7.}

El control verifica la política configurada, los permisos, la consulta de evidencia de prueba dentro de su plazo y la eliminación según su clase. No se prueba una retención de años esperando su transcurso: se inspeccionan configuración, inventario y protección y se ensayan las reglas con datos controlados. La Oferta Económica valoriza los recursos, volúmenes y archivo de estas mismas clases (RT-14.08).
